\SourceLeaf{136}
15, bei dem Zinsfuss von \(3\tfrac12\) proc. für höchstens 17, bei dem Zinsfuss von 4 proc. für höchstens 20 Pensionirte, und der Zinsertrag jener Kapitalvermehrung wird, wenn diese Zahlen erreicht sind, dadurch gänzlich absorbirt. Steigt also die Anzahl der zu pensionirenden resp. über diese Grenzzahlen, so wird zur Bestreitung aller übrigen Pensionen nichts vorhanden sein, als 1) der Ueberschuss, den die jährlichen Einnahmen der Kasse, vor den Erhöhungen von Kapitalien und Pensionssätzen gewährt hatten, oder vielmehr gewährt haben würden, wenn die resp. Normalzahl der Pensionen (15; 17; 20) damals Statt gefunden hätte\footnote{In der Wirklichkeit hatten die Einnahmen damals (1794) schon für 18 Pensionirungen nicht ganz ausgereicht, obgleich zu jener Zeit noch ein jährlicher Zuschuss von 150 Thl. aus der Kirchenkasse geleistet wurde, der später aufgehört hat. Der obige Ueberschlag erleidet aber eine Modification, weil, vorzüglich in Folge von Irregularitäten während der westphälischen Regierung, die Progressionsnorm nicht genau befolgt ist. Hätte man sich ganz strenge daran gehalten, so würden die Pensionssätze seit 1815 immer schon bei geringerer Kapitalhöhe, als geschehen ist, haben erhöhet werden müssen, — und das jetzige Kapitalvermögen würde um vielleicht 5000 Thl. ärmer sein.}. 2) Die etwaige seit jener Zeit bewirkte Steigerung der Apothekenpacht. 3) Der vergrösserte Ertrag der Beiträge der Mitglieder, wegen ihrer gewachsenen Anzahl. Diese precären und schwachen Hülfsquellen würden aber bei weiterm Ueberschreiten jener Normalzahlen bald erschöpft sein, und desto schneller, je höher der Pensionssatz selbst schon angewachsen ist. Diese letztere Bemerkung ist insofern von grosser Wichtigkeit, weil daraus auf das klarste hervorgeht, dass die aus consequenter Befolgung der zweiten Interpretation entspringende Gefahr desto grösser wird, je mehr Pensionserhöhungen bis zum Ueberschreiten der Normalzahl schon Statt gefunden haben.

Wenn übrigens oben bemerkt ist, dass in dem gedruckten Regulativ gar kein Vorbehalt zu finden ist, wodurch einem früh oder spät aus dieser Quelle entspringenden Verbluten der Kasse Einhalt gethan werden könnte, so darf ich nicht verschweigen, dass in die Quittungsformulare, auf welche die Witwen ihre Pensionen erheben, die Bevorwortung aufgenommen ist, dass die Pension auf .\,.\,.\,.\,. Rthl. für jetzt, und \textit{so lange der Kasse Umstände solches gestatten werden}, festgesetzt sei. Wahrscheinlich werden wenige Mitglieder der Witwenkasse diese Clausel kennen: mir selbst wenigstens ist sie, obgleich ich 38 Jahre Theilnehmer gewesen bin, erst ganz vor kurzem bekannt geworden. Sie ist jedoch nicht bestimmt, wie ich anfangs vermuthete, möglichen aus der Progressionsnormirung zu besorgenden Gefahren vorzubeugen; denn sie steht, und zwar genau mit denselben Worten, auch schon in den gedruckten Quittungsformularen vor 1794. Solche in ganz allgemeinen Ausdrücken abgefasste Vorbehalte mögen in einigen Beziehungen ihr Gutes haben: im gewöhnlichen Laufe der Dinge aber bleiben sie, wenn nicht ausserordentliche Veranlassungen eintreten, so lange ohne Anwendung, bis die höchste Noth zwingt. Es ist kein Entscheidungsmerkmal, kein Maassstab angegeben, woran man erkennen kann, ob der Kasse Umstände Zahlung gestatten oder nicht gestatten. Wartet man so lange, bis man schon wiederholt genöthigt ist, zur Bezahlung der Pensionen die Kapitale mit zu verwenden, so hat man sehr wahrscheinlich schon viel zu lange gewartet, und es wird sich dann bewähren, was [S. 128, Z. 12] bemerkt ist. Jedenfalls ist, rücksichtlich der Behauptung des Rufs der Solidität der Kasse, ein grosser Unterschied zwischen einem Zurückgehen des Pensionssatzes in Folge einer ganz bestimmten feststehenden Regel (wie bei der ersten Interpretation); und einer Reduction der Pensionen, wozu die Kasse sich endlich gezwungen sieht, weil sie eben ihre nach dem
